\section{One Church, Two Canonical Communities}

The decisive intervention of 2007 did not erase Saint Stanislaus Parish and replace it with an Institute institution. It brought two endangered continuities into the same church. The old parish still possessed a functioning church building, a campus, accumulated memory, and obligations that its diminished local membership could no longer easily carry. Milwaukee's community attached to the traditional Roman liturgy had worshippers and a stable liturgical purpose but was about to lose the borrowed home it used at Saint Mary Help of Christians in West Allis. Before that tenancy, the community had moved among such places as the Cousins Center and a dance hall. One community had inherited place without enough people; the other had gathered people without a permanent place.

Archbishop Timothy M. Dolan supplied the institutional hinge. The Institute of Christ the King Sovereign Priest announced on October 15, 2007, that he had invited it to serve the Archdiocese of Milwaukee's traditional Latin Mass apostolate. Mass according to the 1962 Roman liturgical books was to begin at Saint Stanislaus on the first Sunday of Advent. The announcement thanked the Franciscan friars who had served the parish. The archbishop also placed the continuing parish in the Institute's pastoral care, and the archdiocesan arrangement was subsequently expressed through Saint Stanislaus Oratory. The current institutional history dates erection of the oratory to 2008.

That sequence requires exact language. Saint Stanislaus Parish and Saint Stanislaus Oratory are distinct canonical entities. Archdiocesan reporting in 2018 described the church, rectory, and school as parish property; the current city assessor lists the historic parcel under St Stanislaus Properties LLC. The oratory is the apostolate organized for worship according to the liturgical books of 1962. The Institute serves both; no checked record says that it bought the historic campus, absorbed the parish, or turned parish property into Institute property. At the study cutoff the same Institute canon is styled parish administrator and rector, two offices held by one priest. A parish Mass in English continues on the Saturday vigil, while the oratory's Sunday and weekday schedule uses the older Roman books. Common clergy, office, bulletin, council work, fundraising, and physical space make cooperation visible, but practical unity does not dissolve juridic distinction.

\statusframe{Archbishop Dolan's 2007 invitation and entrustment; ICKSP institutional history; recurring parish-and-oratory bulletins; archdiocesan parish listing.}{The parish continued, while a distinct oratory received a stable home at the parish church and Institute clergy assumed the linked offices of parish administrator and oratory rector.}{No checked record supports suppression of the parish, transfer of ownership to the Institute, or the claim that ``the parish became an oratory.'' The formal decree itself was not located for this study.}

The arrangement did more than allocate Mass times. It aligned the inherited church with a form of worship that made its deep sanctuary, high altar, communion rail, choir loft, organ, sacristies, processional spaces, and extensive vestment collection useful together. Features that had become expensive remnants could again function as parts of an integrated liturgical environment. This affinity did not make every later design decision historically inevitable, nor did it prove that one moment in the building's past was its only authentic form. It did create a durable reason to research, repair, and use the fabric rather than merely stabilize it.

The people arriving at Saint Stanislaus were not simply neighborhood replacements. The 2018 archdiocesan newspaper found families traveling from as far north as Kohler and as far south as Racine. The revived constituency was therefore regional, assembled by liturgical attachment rather than Polish residence around Mitchell Street. That difference matters. Nineteenth-century Saint Stanislaus concentrated an immigrant population and helped it build a neighborhood. The post-2007 apostolate gathered a dispersed population at an inherited neighborhood landmark. The first pattern sent daughter parishes outward as Polish settlement spread; the second drew households inward across metropolitan and diocesan distances.

Yet the two populations shared a practical resemblance. Early immigrant families assessed themselves, mortgaged their future income, donated bells and marble, and supplied labor because worship required a place. The newer community likewise funded capital work, volunteered for cleaning and construction tasks, recovered furnishings, and organized repeated benefits. The resemblance is not genealogical---most present worshippers cannot be presumed descendants of the founders---but institutional. A stable congregation again treated the building as a common obligation rather than an attractive commodity.

This also clarifies what ``resurrection'' can and cannot mean. Saint Stanislaus was not canonically dead and re-founded. English-language parish worship did not disappear, the Polish foundation was not repudiated, and the oratory did not acquire the parish's history by substitution. Resurrection here is a historical description of recovered capacity: fuller worship, dependable clergy, new families, organized societies, capital maintenance, and an artistic program directed toward use. The old body remained; a new alliance made it act at scale again.

\claimcheck{Entrustment is not ownership}{The phrase ``under the Institute'' is safe only when it describes pastoral care, clerical service, or the oratory apostolate. It is misleading when used to imply that the Institute owns the parish church or that Saint Stanislaus Parish ceased to exist. The continuing parish listing, parish Mass, public property account, and dual title ``parish administrator and rector'' all require the two-community account.}
